% page 430 of the facsimile (image p-429 of the scan)
% block 157.5 x 254.9 mm ; left margin 8.0 mm
\pgc{-12.700mm}{0.000mm}{
\l{}
\l{}
\l{\cel{3}422.}
\l{}
\l{\sou{paranukleino}.\cel{2}(\sou{kemio biol}.) Substanco ye qua konsistas la}
\l{\cel{3}nukleoli di la celulo-nukleo. -}
\l{}
\l{\sou{parapeto}.\cel{2}I. Muro, ye la nivelo di la kudo, alonge la bordo}
\l{\cel{3}di ponto, kayo, teras-tekto, por impedar la neprudenti}
\l{\cel{3}falar aden precipiso, abismo, maro, e c. - II. Masivo ek}
\l{\cel{3}tero o masonuro qua quaze kronizas remparo, e shirmas kon-}
\l{\cel{3}tre la pafi de la enemiko. - DEFIRS.}
\l{}
\l{\sou{paraplasmo}. (\sou{kemio biol.}) Substanco vitratra, mi-fluida, qua}
\l{\cel{3}konsistas ye globulino, ed en qua imersesas la elementi di}
\l{\cel{3}citoplasmo.- DEFIS.}
\l{}
\l{\sou{parazito}. I. (en\cel{1}\sou{la epoki antiqua}). Repasto-gasto singladia di}
\l{\cel{3}richo, qua havis kom tasko amuzar ica. - II. Ta qua mes-}
\l{\cel{3}tiere su invitigas por repastar. - III. Planto od insekto}
\l{\cel{3}qua vivas ek altra animalo od altra planto. - DEFIRS.}
\l{}
\l{\sou{parcelo}. Ek ensemblo, mikra parto di qua la formo esas deter-}
\l{\cel{3}minita precize. - DEFIS.}
\l{}
\l{\sou{pardonar}. (\sou{trans.})\cel{2}I. (\sou{ulo ad ulu - o : ulu pri ulo)}\cel{1}Remisar}
\l{\cel{2}(kulpo.)\cel{2}II. (en\cel{1}\sou{la formulo di politeso)}. Indulgar.- EFIS.}
\l{}
\l{\sou{parear}.\cel{2}(\sou{trans.}) Evitar atingeso, per deviacigar (la frapo-}
\l{\cel{3}stroko). - DEFIS.}
\l{}
\l{\sou{parenkimo}. (\sou{anat.})\cel{2}I. Tisuo sponjatra, qua karakterizas la}
\l{\cel{3}viceri. - II. (\sou{bot.})\cel{2}Tisuo utrikula di la planti. - DEFIRS}
\l{}
\l{\sou{parento}.\cel{2}Persono qua esas membro di la familio (di ulu)}
\l{\cel{3}per nasko o mariajo. - EFIS.}
\l{}
\l{\sou{parentezo}. I. (\sou{gram.)}\cel{3}Frazo acesora, komplementa, e, kom}
\l{\cel{3}tala, provizita ye signo partikulara qua izolas lu interne}
\l{\cel{3}di la frazo precipua. - II. Singla de la signi :\cel{2}(), qui}
\l{\cel{3}indikas parentezo. - DEFIS.}
\l{}
\l{\sou{parfumo}. I. Odoro agreabla,naturala. - II. Substanco qua igas}
\l{\cel{3}odorar tale. - DEFIRS.}
\l{}
\l{\sou{parhelio}. (\sou{astron.)}\cel{4}Fenomeno lunala qua produktesas da la}
\l{\cel{3}reflekteso di lumo sur la kristali glacia qui suspendesas}
\l{\cel{3}en atmosfero. - DEFIS.}
\l{}
\l{\sou{pariar}.\cel{2}(\sou{trans., pri, por, kontre}).\cel{2}Konvencionar pri pe-}
\l{\cel{3}kunio-quanto riskenda (generale, da singla koncernato)}
\l{\cel{3}por ta qua esabos justa pri kozo kontestata. - DFR.}
\l{}
\l{\sou{\textquotedbl{}paria}\textquotedbl{}.En India, persono qua apartenas a la kasto maxim in-}
\l{\cel{3}fra e qua, kom tala, esas la objekto di desestimo.}
\l{}
\l{\sou{pari}eto. I. Separilo ek masonuro lejera, menuzuro, quan onu}
\l{\cel{3}erektas en domo, apartamento, e c. - II. To quo dividas}
\l{\cel{3}kavajo aden du o plusa parti. - EFIS.}
}
